\documentclass[11pt]{article}
\usepackage[T1]{fontenc}
\usepackage{lmodern}
\pdfmapfile{+lm.map}
\pdfmapfile{+cm.map}
\pdfmapfile{+cmextra.map}
\pdfmapfile{+symbols.map}
\usepackage[a4paper,margin=28mm]{geometry}
\usepackage{amsmath,amssymb,amsthm,booktabs,longtable,array}
\usepackage[protrusion=true,expansion=false]{microtype}
\usepackage{xurl}
\usepackage[hidelinks]{hyperref}
\hypersetup{pdftitle={Compact Fixed Controller S Computation from a 38 Phase Cyclic Tag System},pdfsubject={A formal compact source construction with explicit selected path semantics},pdfkeywords={S combinator,cyclic tag systems,formal verification,Lean}}
\usepackage{enumitem}
\setlist{itemsep=2pt,topsep=5pt}
\setlength{\emergencystretch}{2em}
\lefthyphenmin=2
\righthyphenmin=3
\newcommand{\code}[1]{\begingroup\normalfont\small\upshape\def\UrlBreaks{\do\.\do\_}\def\UrlBigBreaks{\do\/}\nolinkurl{#1}\endgroup}
\newcommand{\NN}{\mathbb{N}}
\newcommand{\Bool}{\mathbb{B}}
\newcommand{\Terms}{\mathcal{T}_{S}}
\newcommand{\StepS}{\mathrel{\longrightarrow_{S}}}
\newcommand{\enc}{\mathsf{enc}}
\newcommand{\obs}{\mathsf{obs}}
\newcommand{\readout}{\mathsf{read}}
\newtheorem{theorem}{Theorem}[section]
\newtheorem{lemma}[theorem]{Lemma}
\newtheorem{proposition}[theorem]{Proposition}
\newtheorem{corollary}[theorem]{Corollary}
\theoremstyle{definition}
\newtheorem{definition}[theorem]{Definition}
\newtheorem{remark}[theorem]{Remark}
\title{Compact Fixed Controller S Computation\\from a 38 Phase Cyclic Tag System}
\author{}
\date{7 October 2026}

\begin{document}
\begin{center}
{\LARGE Compact Fixed Controller S Computation\\[4pt]
from a 38 Phase Cyclic Tag System\par}
\vspace{12pt}
7 October 2026
\end{center}
\vspace{6pt}
\begin{abstract}
A fixed finite controller can select native S-combinator contractions that
simulate a compact universal cyclic tag system. We give a concrete
38-phase instance, an explicit compiler from finite natural-register machines,
and a direct compiler from conventional integer-tape Turing machines with
arbitrary finite alphabets. A fixed regular tree language recognizes source
halting. At its first accepted sample, a fixed current-tree decoder returns
the exact register-machine result. The controller restarts at the root on
each invocation, retains no state between invocations, and has an all-input
linear invocation bound. The generic cyclic-tag-to-S construction is due to
Cinematic Strawberry, whose earlier 912-phase construction already establishes
fixed-controller universality for Boolean tapes. The contributions here are
the compact source route, its explicit formal compiler, and the generated-origin
and first-event numerical composition for this instance. All 67 project modules
passed fresh Lean 4.33.1 kernel replay. Complete operation-cost interpretations
are stated separately from the checked semantic and size lemmas.
\end{abstract}

\section{Introduction}

The sole rewrite rule of the S combinator is
\[
       S\,x\,y\,z \;\StepS\; x\,z\,(y\,z),
\]
with application associated to the left. The computational question depends
on the encoding, reduction strategy and observation interface. This paper
studies a selected trajectory of finite closed S trees, with source halting
recognized during continuing reduction. The trajectory is generated by one
fixed finite root-reset controller. All auxiliary interfaces are specified
before the source machine and input are supplied.

Cinematic Strawberry's prior construction~\cite{strawberry} establishes this
kind of computation through a fixed 912-phase cyclic tag system. Its public
headline already includes regular-language halting on the actual root-reset
trajectory, conventional Boolean infinite-tape semantics, and exact source-row
and terminal-bit readout~\cite{headline,infinite,prioroutput}. We use its generic
cyclic-tag-to-S encoding, scheduler and finite selector. The present result is
a different compact instantiation and an explicit proof of its source,
intermediate-event and numerical-output obligations.

The source endpoint is ais523's 19-symbol alternating tag system UT19~\cite{ut19}.
A one-hot translation gives a 38-phase binary cyclic tag system. An explicit
compiler supplies a universal finite-register source with an exact numerical
output convention. A further checked compiler begins with a conventional
integer-indexed tape and an arbitrary finite alphabet. The main proof tracks
all intermediate S samples: a completed phase-17 pattern can be copied into
several parts of a tree, so checkpoint correspondence alone does not establish
that the first matching occurrence carries the intended result. Generated
label and audit invariants prove that every match in the first accepted tree
has the same frozen payload.

The contribution is therefore threefold. First, the fixed cyclic tag endpoint
has period 38 with an explicit machine/input compiler. Second, the first-event
result theorem covers every finite natural-register program and input, with
no residual source-event, liveness or witness-correctness hypothesis. Third,
the composition has a reproducible kernel-checked implementation, including a
direct finite-alphabet tape bridge. The change from 912 to 38 is a change in
one construction parameter; it implies no corresponding factor of improvement
in runtime, control-state count or term size.

\section{The computational model}

\begin{definition}[Terms and native steps]
The set $\Terms$ consists of finite terms generated by $S$ and binary
application. Its occurrence size is
$|S|=1$ and $|uv|=1+|u|+|v|$. A native step contracts exactly one occurrence
of $Sxyz$ to $xz(yz)$ inside an arbitrary application context.
\end{definition}

Sizes count occurrences, including every repeated subtree. A compressed DAG
implementation may use fewer distinct nodes; all asymptotic readout statements
in this paper use occurrence-tree size.

\begin{definition}[Root-reset finite selector]
A controller has a finite control set and an occurrence cursor on the current
tree. A cursor is a focus with its parent-frame path. A transition observes
only the current control, node kind and incoming side. Primitive commands
change control, move along one child or parent edge, contract a focused native
S redex, or reject. Each invocation starts in the same control at the root.
Before the selected contraction, navigation is read-only. Between invocations
only the resulting bare S term is retained.
\end{definition}

The cursor may have unbounded depth as the input tree grows; finiteness concerns
control, not tree size. The imported selector contract makes this distinction
explicit. Its inter-invocation state is literally \code{Unit}. For the fixed
controller $C$ used below there is a constant $a$ such that every invocation on
an arbitrary finite tree $T$ stops within $a(|T|+1)$ primitive microticks. On a
normal tree it makes no contraction. On every encoded trajectory in the theorem
it selects exactly one native contraction and returns its successor. The
constant $a$, controller and observation algorithms are independent of the
source machine and input. The bound is per invocation, not per simulated
computation.

Write $F_C:\Terms\to\operatorname{Option}(\Terms)$ for the selected successor.
For an encoded source instance $z$, define $T_z(0)=E(z)$ and
$F_C(T_z(i))=\operatorname{some}(T_z(i+1))$. The proof establishes this equation
for every $i$, so the totalization used in the implementation is immaterial
on encoded paths. Observation occurs at each resulting tree, including the
initial tree. For a fixed Boolean observer $H$, define
\[
 \operatorname{First}_H(T,n)\iff H(T(n))\ \land\
        \forall j<n,\ \neg H(T(j)).
\]
The decoder $D:\Terms\to\operatorname{Option}(\NN)$ receives only the current
tree. It has no source program, input, trace, checkpoint index or execution
oracle as a runtime argument.

\subsection{Source machines}
A register machine has finitely many labels and $R\geq1$ natural registers.
Its instructions are increment and jump, decrement with separately specified positive
and zero successors, and halt. A zero test leaves a zero register unchanged;
a positive decrement subtracts one. Halting configurations are absorbing.
For input vector $u$, let $M[u]\Downarrow v$ mean that a finite run reaches
halt with register 0 equal to $v$.

The tape source has finite state and alphabet sets, a start state, a blank
symbol, and a finite partial transition table. A transition writes at the
current integer coordinate and moves by $-1$, $0$ or $1$. No transition means
halt. An input word occupies coordinates $0,1,\ldots$, the head starts at $0$,
and the rest of the tape is blank. States contain a genuine total tape
function $\mathbb Z\to\Sigma$. There is no bounded visited-window hypothesis.
Finite transition functions are supplied by explicit finite tables when
measuring the encoding's binary description length.

\section{Uniform computation theorems}

\begin{theorem}[Register results]\label{thm:register}
There are a fixed root-reset finite controller $C$, a fixed regular tree
language $H$, an effective finite-tree encoder $E_R$, and a fixed total decoder
$D$ such that, for every finite register machine $M$, natural input vector $u$
and $v\in\NN$,
\[
 M[u]\Downarrow v\quad\Longleftrightarrow\quad
 \exists n\;\bigl(\operatorname{First}_H(T_{M,u},n)
           \land D(T_{M,u}(n))=\operatorname{some}(v)\bigr).
\]
Every adjacent pair in $T_{M,u}$ is one native S step. Every invocation of $C$
on every finite input tree satisfies the fixed linear bound above. The encoder
has a singly exponential bound in an explicit binary source description; the
decoder has a polynomial all-input bound in current occurrence-tree size.
\end{theorem}

The semantic equivalence is
\code{SOnlyUniversality.result_iff_first_event}. The encoder and decoder are
ordinary executable definitions. The construction of $E_R$ has no halting-time
parameter and never runs the source to determine its initial S term. The
resource interpretation is made precise in Section~\ref{sec:resources}.

\begin{theorem}[Conventional tape halting]\label{thm:tape}
There is an effective encoder $E_T$ for every finite-alphabet tape machine $M$
and finite word $w$, using the same $C$, $H$ and $D$, such that
\[
 M\text{ halts on }w\quad\Longleftrightarrow\quad
       \exists n\ H(T_{M,w}(n)).
\]
Equivalently, a first accepted sample exists. If the source does not halt,
$H$ rejects every finite sample. At the first accepted sample, $D$ returns the
code of the finite left stack in the specified two-stack simulation.
\end{theorem}

The halting endpoint is \code{SOnlyTuringUniversality.halting_iff_event}; the
value-specific endpoint is
\code{SOnlyTuringUniversality.left_stack_result_iff_first_event}.
The left stack retains explicit blank cells and visited extent. Its code is
an observable of this deterministic simulation history, rather than a
canonical arbitrary final tape-output word. The general natural-result
interface is supplied by Theorem~\ref{thm:register}.

\section{Finite registers and the numerical BP2 frontend}
\label{sec:source-registers}

The source reduction is an executable composition, with no supplied bound on
source running time.  Its intermediate mechanisms follow ais523's
Amnesiac From Minsk~\cite{amnesiac},
Waterfall Model~\cite{waterfall},
and Brainpocalypse II~\cite{bp2}
constructions; the Waterfall account also credits zzo38's independent model.
The explicit register-frontend tables, compiler composition, and proofs
described here are this development's contributions.  The generic CTS-to-S
encoding and root-reset controller are supplied by
\href{https://github.com/cstrawberry/predictive-universe/tree/85a867988442fc423279341200f81634a1e65582/docs/paper/related/pure_s_universality}{Cinematic Strawberry's pinned construction}.

A source machine has $R\geq1$ natural registers, $Q\geq1$ labels, an entry
label, and instructions $\operatorname{INC}(i,t)$,
$\operatorname{DECJZ}(i,t_+,t_0)$, and $\operatorname{HALT}$.
Zero testing leaves a zero register unchanged.  BP2 instead executes a finite
list of $+i,-i$: positive decrements advance normally, while a zero decrement
sets that counter to one and restarts at command zero, retaining all other
counters.  Falling off the list halts; all BP2 counters initially equal zero.

\begin{theorem}[Exact register frontend]
For every such $M$ and input tuple $a$, the total function
\code{SOnlyMachineNat.compile} produces a nonempty finite BP2 program $B(M,a)$
with dense counter labels.  It halts exactly when $M(a)$ halts.  For each
$v\in\NN$, it halts with first counter $2v+3$ exactly when $M(a)$ halts with
register zero equal to $v$.
\end{theorem}

\begin{proof}[Construction and proof]
First use positive amnesiac counters: $+i$ increments and follows a fixed
successor $p_i$; $-i$ decrements above one and follows $s_i$, or leaves one
unchanged and follows $f_i$.  Lean stores these values minus one.  Keep
$D_i=r_i+1$ and allocate five helpers $I_q,S_q,F_q,U_q,V_q$ at \emph{every}
label, initially one.  Let $E(q)$ be $+I_q$, $+S_q$, or halt according to the
source instruction.  The successor table is
\[
\begin{array}{c|ccc}
 &p&s&f\\\hline
D_i&-I_0&-S_0&-F_0\\
I_q&+D_i&E(t)&-I_{q+1}\\
S_q&+F_q&-V_q&-S_{q+1}\\
F_q&-D_i&-U_q&-F_{q+1}\\
U_q&-F_0&E(t_+)&+V_q\\
V_q&-S_0&E(t_0)&+U_q
\end{array}
\]
Here instruction-dependent entries apply to the appropriate instruction;
inapplicable instruction-dependent entries and end-of-chain fallbacks are halt.  Scans visit all
labels.  A chain with exactly one marked value two passes the ones, restores
the mark to one, and takes precisely its success edge.  An increment marks
$I_q$, increments $D_i$, and scans $I$.  A decrement marks $S_q,F_q$, tests
$D_i$, then uses the $S,V,U,F$ or symmetric $F,U,V,S$ route to restore every
helper and select $t_+$ or $t_0$.  Thus the instruction-boundary invariant is
exact, including shared data registers and repeated call sites
(\code{SOnlyMachine.macro_simulation}).

For $k=R+5Q$ amnesiac counters, introduce $4k+1$ Waterfall clocks: data $X_i$,
action clocks $Q_a$ for $a=+i,-i,H$, and recovery clocks $T_i$.
At normalized action $(c,a)$ set $X_i=2c_i$, $Q_a=0$, and every other control
clock to four.  With $e_Z$ a unit vector, the signed trigger rows are
\[
\begin{aligned}
b_{Q_{+i}}&=2e_{X_i}+4e_{Q_{+i}}-4e_{Q_{p_i}},&
b_{Q_{-i}}&=-2e_{X_i}+4e_{Q_{-i}}-3e_{T_i},\\
b_{X_i}&=2e_{X_i}+3e_{T_i}-4e_{Q_{f_i}},&
b_{T_i}&=\mathbf1+3e_{T_i}-4e_{Q_{s_i}}.
\end{aligned}
\]
Use $A_Z=b_Z+5\cdot\mathbf1$ and $A_{Q_H}=0$, adding coincident terms.
All nonhalt entries lie in $[1,9]$.  Subtracting the common minimum after a
trigger cancels its uniform offset.  An increment reaches the next normalized
action; a decrement makes $X_i=2c_i-2$ and $T_i=1$.  If $c_i=1$, $X_i$ is
uniquely minimal and restores the failure state; otherwise $T_i$ is uniquely
minimal and its recovery row restores the successful-decrement state.
These arithmetic identities prove tie-free simulation
(\code{SOnlyMachineClock.step_simulation}).

BP2 implements each clock with a value cell $v_i$ and pending marker $z_i$,
plus one marker $m$.  For a nonnegative vector $b$, write $\operatorname{Add}(b)$
and $\operatorname{Sub}(b)$ for literal unary increments and decrements.
The initializer is
\[
 \operatorname{Add}(a);-m;+m;\operatorname{Sub}(a),
 \qquad a(v_i)=w_i,\quad a(z_i)=1,\quad a(m)=0.
\]
Its first zero guard preserves initialization across a restart; subsequent
visits cancel exactly, even when working cells are zero.
Trigger $i$ uses the same pattern guarded by $z_i$, adding
$A_i-e_i$ to value cells.  A sweep subtracts one from every $v_i$; tests
$-z_i;-v_i;+v_i;+z_i$ mark the unique zero and restart.  The selected trigger
then compensates for the test's automatic increment and applies exactly $A_i$.
The terminal pair $-m;-m$ restarts ordinary sweeps.  The halt trigger instead
adds $2\cdot\mathbf1-e_h$ and increments $m$; one final sweep falls off with
$v=y+\mathbf1$, every $z_i=1$, and $m=0$.  Hence the first value is
$2(r_0+1)+1=2r_0+3$.
Finite positive macro durations and absorbing halt states give reflection at
arbitrary target times, not only checkpoints.  Dense renumbering preserves
the output coordinate.  These are
\code{SOnlyMachineNat.halt_iff} and
\code{SOnlyMachineCorrectness.output_iff}.
\end{proof}

The formal compiler initializes clocks at $V+5$, whereas the earlier selective
Python construction uses $V+1$ and allocates helpers only where needed.
The formal count is $c=8k+3$, $k=R+5Q$; its checked bound, with
$N_c=4k+1$ and $s=\sum_i a_i$, is
\[
 L\leq2c(2s+9)+N_c(18c+2)+c+4N_c+4
 \quad\text{(\code{SOnlyMachineBounds.compile_length_bound}).}
\]
The two variants preserve semantics but are not byte-for-byte identical.

\section{Programmed UT19 execution and its first event}
\label{sec:source-ut19}

The next reduction uses ais523's CC0
UT19 construction~\cite{ut19}.
Its literal productions, in published one-based notation, are
\[
\begin{array}{r@{\,\mapsto\,}l r@{\,\mapsto\,}l r@{\,\mapsto\,}l}
1&(2,3)&2&(4,4)&3&(18,4)\\
4&(1,1,19)&5&(7,9)&6&(8,9)\\
7&(10,10)&8&(11,10)&9&(18,10)\\
10&(5,6)&11&(6,5,19)&12&(14,14,14,14)\\
13&(15)&14&(16,16)&15&(16,17)\\
16&(12,13)&17&(12,12,12,12)&18&(18)\\
19&\epsilon&\multicolumn{2}{c}{}&\multicolumn{2}{c}{}
\end{array}
\]
Word powers denote literal repetition. A microstep removes one symbol, appends its production only in \emph{take}
phase, and toggles take/skip.  The designated event is a \emph{prestate} with
take phase and head 18; a skipped 18 is not an event.

For a BP2 program of length $L$ and $c$ counters, add a halt counter $H=c+1$.
Choose the least power of two $N\geq L+3$ and put $W=2N$.
The seed is
\[
19\,M_1K_1\cdots M_cK_cM_HK_HM_{H+1},
\qquad K_i=(12,13)^2.
\]
Each memory has $W$ cells.  The expanded program is dummy, the $L$ commands,
halt-A, halt-B, and padding to $N$ slots.  Every slot uses two halfcommands;
each halfcommand comprises Command, Parity, and Reset passes, where a pass
consumes exactly its entrance word.

\begin{lemma}[Finite programmed-memory window]
For any prescribed $W$ width parities there is an explicit initializer giving
those observations at times $0,\ldots,W-1$, independently of the incoming
Command alignment history.  Odd Reset restores this initializer exactly.
\end{lemma}

\begin{proof}[Proof sketch]
A cell with dynamic bit $h_i$ is $(5,6)$ or $(6,5,19)$; a fixed inverter bit
$a_i=1$ adds $(1,1,19)$.  Put $c_i=h_i\oplus a_i$.
Direct production calculation gives, under normal Parity and Reset,
$c'=Pc\oplus g\mathbf1$, where $P$ is inclusive prefix XOR and $g=1$ means
skip alignment.  Width is observed \emph{before} updating:
$v_t=\bigoplus_i c_{t,i}$.
For $W=2^d$, the unforced observation transform is
\[
 v_t=\bigoplus_{i:\,(i\mathbin{\&}t)=t}a_i,
 \qquad
 a_i=\bigoplus_{j:\,(j\mathbin{\&}i)=i}v_j.
\]
The superset transform is its own inverse over $\mathbb F_2$.
The formal proof uses even/odd index decomposition and identifies its recursive
transform with the increasing-stride butterfly.  The extensionally equal increasing-stride implementation uses
$d=\log_2W$ stages and $W\log_2W/2$ bit XORs. This count concerns that
implementation; the output-equivalence theorem is not an exact operation-count
theorem for direct evaluation of the recursive Lean definition.  An alignment disturbance in
update $t$ first affects observation $t+W$; reset sets every $h_i$ to zero.
These facts are \code{SOnlyTemporalMemory.forced_initializer_correct},
\code{initializer_eq_increasing_strides}, and
\code{SOnlyMemory.memory_reset_vector}.  The emitted choice
$d=\lfloor\log_2(L+2)\rfloor+2$ satisfies
$2L+3<W\leq4L+8$ by \code{SOnlySimulation.compiler_capacity} and
\code{compiler_width_bound}.
\end{proof}

Memory widths are adjacent XOR differences of the desired component
alignments, so prefix XOR telescopes to exactly those alignments.
\code{SOnlySimulation.scheduledWidth_eq_published} checks equality with the
literal width recipes, including unused halt-B and padding slots.
For $\mathcal C(n)=(12,13)^{2^n}$, ordinary BP2 value $x$ is
$\mathcal C(2x+1)$.  A halfcommand increments $n$ when Command starts take;
a positive-exponent decrement subtracts one under even Reset.
Thus $+i$ increments all exponents, then increments $i$ and decrements the
others.  Similarly, $-i$ decrements all, then decrements only $i$ and
increments the others.  These implement changes of two in the target
exponent and zero elsewhere.

The zero branch requires more than ordinary decrement arithmetic.
Decrementing $\mathcal C(0)$ produces a singleton 15 in Parity.  At a source
zero decrement it is selected, creates odd Reset, and leaves the protected
word $D=(12,12,12,12)$.  Every other counter is incrementing in that
halfcommand: no positive decrement is exposed to odd Reset, which would
otherwise produce the wrong protected form.  All memories reset exactly;
any 18 they emit is skipped.  The restarted dummy decrements all and then
increments all, but $D$ responds to the first half by becoming
$\mathcal C(2)$, then $\mathcal C(3)$.  The target therefore becomes one,
as BP2 requires, while other values are retained.  The initial leading 19
supplies the same skip alignment for the first dummy.  This proves closure
through arbitrarily many restarts
(\code{SOnlySimulation.source_zero_decrement_run}).

At command $i$, memory time is $2i+2$.  Each epoch either resets or reaches
halt-A's last needed observation $2L+3<W$, so the memory lemma is never used
past its protected window.  Halt-A decrements all exponents, then increments
ordinary counters and decrements the halt counter with odd Parity.  The latter's
singleton 15 is skipped and disappears; Reset begins take with head 18.
Exactly five passes after halt-A entrance, the event queue is
\[
 R_1\,16^{4^{x_1+1}}\cdots R_c\,16^{4^{x_c+1}}R_HT,
\]
where all separators are nonempty and contain no 16.  The $x_i$ are the final
BP2 values.  Thus maximal 16-runs give precisely the retained counters
(\code{SOnlySimulation.finalEvent_readout}); repeated division by four
recovers them.  The first counter then yields $(x_1-3)/2$ for register output
(\code{finalEvent_readMachineValue}).

Each preceding simulated step is a finite nonempty, event-free run.
Unconsumed entrance symbols prevent exhaustion inside a pass, and memories
supply nonempty output at its end.  Positive simulation durations make
nonhalting checkpoints cofinal, excluding an event at \emph{every} finite
microstep.  Consequently \code{SOnlySimulation.source_halting_iff_event}
is a two-way unbounded theorem, and \code{first_event_exact} identifies the
actual first event queue, rather than only an expected macro-boundary.

\section{The fixed 38-phase interface and conventional tapes}
\label{sec:source-cts}

Encode symbol $s$ by $\eta(s)=0^{s-1}1\,0^{19-s}$.
The fixed CTS has the nineteen encoded productions followed by nineteen empty
appendants.  Take and skip boundaries are phases zero and nineteen.

\begin{lemma}[Exact block and event transfer]
Nineteen CTS steps simulate one UT19 microstep.  At every CTS time $t$, the
phase-17/head-one event holds exactly when UT19 is selected at microstep
$\lfloor t/19\rfloor$ and $t\bmod19=17$.  A first UT19 event at $n$ therefore
becomes the first CTS event at $19n+17$.
\end{lemma}

\begin{proof}[Proof sketch]
After consuming a prefix, FIFO order leaves its unconsumed suffix ahead of
all generated appendants.  Scanning one 19-bit block appends exactly its
production in the take half, or nothing in the skip half, and toggles the
boundary phase.  The remaining original bits prevent premature exhaustion.
The only one in $\eta(s)$ is at offset $s-1$; phase 17 can therefore signal only
selected symbol 18.  Division with remainder covers arbitrary times.
These are \code{SOnlySource.trajectory_simulation},
\code{event_at_arbitrary_time}, and \code{first_event_exact}.
\end{proof}

At the event, the current CTS word is $10\,\eta(\text{tail})$, not a block-aligned
encoding.  Restoring the known seventeen-zero prefix reconstructs the whole
UT19 event word; \code{SOnlyResultBits.decodeEventWord_correct} proves this
structural readout.  Composition gives
\code{SOnlySimulation.compiled_halting_iff_cts_event} and
\code{compiled_noPrematureEmpty}, the exact source interfaces consumed by the
fixed S backend.  In particular, UT19's rule $18\mapsto18$ is an event marker,
not a claim that the tag queue or S term normalizes.

Finally, the source need not be assumed universal by convention.
\code{SOnlyTapeStacks} treats a deterministic finite-alphabet machine on an
integer-indexed tape, with left, stay, and right moves and absent transition
meaning halt.  Finite control stores the state and scanned symbol; nearest-first
left and right stacks represent every remaining tape cell, with empty pop
returning blank.  For alphabet size $A$, use base $B=A+1$ and
\[
 \operatorname{code}([])=0,\qquad
 \operatorname{code}(a::u)=(a+1)+B\operatorname{code}(u).
\]
Nonzero digits preserve explicit blanks and make this code injective.
Literal INC/DECJZ macros implement push by multiplication and addition, and
pop by a finite-control residue cycle and quotient transfer.  Registers zero
and one hold the stacks; register two is scratch, restored to zero.
The assembled table has
$|Q_{\rm TM}|A(5(A+1)+5)+1$ labels and depends only on the machine; its entry
selects the input's first symbol and the right register contains the remaining
input.  Positive-duration macro execution and cofinal checkpoints reflect
halting even inside target macros
(\code{SOnlyTapeCompiler.input_halting_iff}).

Hence \code{SOnlyTuringUniversality.halting_iff_event} covers arbitrary finite
inputs without a tape-window or runtime premise.  Arbitrary register programs
retain their exact natural output.  For this TM bridge the output is specifically
the finite left-stack code, including retained blanks and visited extent:
\code{left_stack_result_iff_first_event}.  It is a history-dependent observable
of this simulation, not a canonical final tape-output word.


\section{The fixed S implementation}\label{sec:backend}

Let $P$ be the literal 38-phase cyclic tag program obtained above. Its true-bit
appendants contain 760 bits in total, of which 40 are ones. Fix its canonical
balanced action dispatcher $A$. There are 76 phase/bit leaves. The leaf for
phase 17 and bit true has route $0100011$, with left and right encoded by 0
and 1. Define the S generator and selector by the pinned upstream operations:
\begin{align*}
 E_P(b)&=\operatorname{generator}(\operatorname{compileActions}(P,A),b),\\
 F_P&=\operatorname{RootResetFiniteAllInputsTraceAgreement.selector}(P).
\end{align*}
Here $b$ is the compiled initial bitword. $P$, $A$, the compiled action code and
$F_P$ are fixed. Source-specific information is confined to $b$.

\begin{proposition}[Imported generic realization]\label{prop:generic}
For every finite bitword $b$ and every sample index $i$, the root-reset selector
returns the next term of the generic persistent CTS simulation. That successor
is obtained by one native S contraction. Public checkpoints decode to the
corresponding iterates of $P$. The actual finite controller realizes this same
selector on every input tree and satisfies its all-input stopping bound.
\end{proposition}

\begin{proof}
Use the following pinned generic results, specialized at $P$:
\begin{itemize}[nosep]
\item \code{RootResetFiniteAllInputsTraceAgreement.agreesOnEverySample};
\item \code{termOnlyPath_eq_persistentPath};
\item \code{RootResetFinitePrioritySelector.selectorContract}.
\end{itemize}
The local specialization \code{SOnly38.selector_eq_projected} identifies the
bare-term selector with the invocation contract, and
\code{SOnly38.trajectory_selects} and \code{trajectory_contracts} identify every
sample, not just checkpoints. The generic dependency closure is attributed
to Cinematic Strawberry~\cite{strawberry,generic} and replayed without replacing
its construction by a Python implementation.
\end{proof}

\subsection{A fixed completed-response pattern}
The designated observation is a completed, fresh Local response shell for
label $(17,\mathsf{true})$. The shell contains the fixed dispatch route and
nineteen completed appender-history arguments. Call its independent-hole
pattern $F_{17}$. The language $H$ consists of finite S trees containing an
occurrence matching $F_{17}$. Holes place no equality constraint on distinct
subtrees. Equal-variable matching is not silently used to obtain regularity.

\begin{lemma}[Finite observer]\label{lem:regular}
$H$ is recognized by a fixed bottom-up tree automaton on every finite tree.
\end{lemma}
\begin{proof}
Store one Boolean match bit for each occurrence in the fixed pattern, and one
bit recording a match at any descendant. At an S leaf, a hole and an S pattern
match, while an application pattern fails. At an application node, the bit
for an application subpattern is the conjunction of its left and right child
match bits. A hole always matches. The descendant bit is the disjunction of
the root match and the child descendant bits. Structural induction proves
agreement with pattern matching and descendant occurrence. The complete
Boolean cube has $2^{|F_{17}|+1}$ states, including unreachable states, giving
an explicit finite cover. This is checked in \code{SOnlyObserver}, in particular
\code{event_finite_cover} and its transition-agreement lemmas.
\end{proof}

The size of this uncompressed cover is a fixed construction constant, not a
claim of an efficient minimal automaton. A separate implementation can share
equal subpatterns; the theorem uses the occurrence-based cover above.

\section{Global first-event provenance}\label{sec:first}

Generic checkpoint simulation does not by itself identify the earliest
occurrence of the newly chosen pattern. A native S contraction copies a
subtree. Completed shells may survive in inactive fields or in frozen audits,
and a response is assembled over many contractions. Correctness therefore
requires an invariant covering every relevant descendant at every intermediate
sample.

\subsection{Licenses for generated shells}
A proof-only license is an allowed pair of source label and frozen snapshot
value for a fresh completed Local. The invariant is extensional: it does not
assign an address-indexed causal identity to each physical copy. Licenses are
not runtime data supplied to the selector or decoder. The preservation proof follows actual constructors:
canonical deletion, frame and route steps, appender responses, rebuilt Local
ancestors, clock/fuel phases, jobs and complete scheduler stages. Copying an
already licensed subtree preserves its license; a newly completed response
receives the label of the consumed CTS prestate and the literal post-deletion
carrier as its audit. Registered ancestor reconstruction preserves the fields
that identify these licenses.

\begin{lemma}[Generated origin invariant]\label{lem:origins}
In every generated scheduler prefix whose source CTS prestates are nonempty,
each fresh completed Local occurrence has a licensed label/snapshot pair
from the actually executed source-response prefix. Its frozen audit equals
the associated literal deleted carrier. No selected label can first appear before a response
with that source label completes.
\end{lemma}

\begin{proof}[Proof structure]
The local induction is proved for the actual mutation constructors, rather
than an abstract closure assumption on all possible terms.
\code{SOnlyProvenance}, \code{SOnlyGeneratedOrigins} and
\code{SOnlyCanonicalOrigins} establish local preservation, including deletion
and copied fields. \code{SOnlySchedulerOrigins} and the global ancestor,
terminal and stage modules lift this invariant through the concrete scheduler.
Exact mutation chains list post-contraction configurations in order. Induction
on their constructors equates list position $j$ with the corresponding global
sample, and concatenation preserves that chronology. The resulting global
invariant has no caller-supplied origin annotation in the final theorem.
\end{proof}

A stage of horizon $h$ runs jobs that restart from the original bitword and
execute source responses $0,\ldots,h-1$ in order. This repeated-prefix schedule
is part of the imported construction. Source nonemptiness up to the first
designated event, established by the compiler, is sufficient for the first-event
argument. For nonhalting compiled inputs the source remains nonempty forever.

\subsection{Sharp creation time and a common audit}
Let $r$ be the first CTS prestate whose phase is 17 and whose head bit is true.
Earlier executed responses have different labels. Thus Lemma~\ref{lem:origins}
excludes an $F_{17}$ match in every earlier job and in the prefix before the
first response to $r$. The exact response trace then resolves the remaining
interior samples.

\begin{lemma}[First acceptance and witness agreement]\label{lem:first}
For a compiled bitword $b$ with first designated CTS event $r$, the first
accepted S sample is
\[
 \tau(b,r)=\operatorname{firstJobPreResponseTime}(b,r)+57.
\]
Every occurrence matching $F_{17}$ in that tree has the same frozen audit at
relative address $\mathsf{LLLR}$. That audit publicly decodes to the tail of
the consumed CTS word.
\end{lemma}

\begin{proof}[Proof structure]
One selected C4 contraction precedes the 56-contraction response. In the
response's earlier positions, the new shell is not yet the completed pattern;
previous licenses still exclude the designated label. Its final contraction
creates the first complete shell with that label. The response-position
lemmas are lifted to the whole tree, including copies and reconstructed
ancestors. They show that every match in the first accepted sample has the
same creation snapshot, not merely that one convenient occurrence is correct.
The public seed-free carrier decoder agrees with the semantic decoder on this
snapshot. The checked result is
\code{SOnlyGlobalFirstEvent.firstSourceEvent_current_read}, with structural
witness agreement supplied by the first-event and global-origin modules.
\end{proof}

The offset 57 is a native-contraction offset in this exact pinned response,
not a general bound on simulation time. Scheduler time before the response
can be much larger.

\subsection{A decoder on the current tree}
The fixed decoder searches in root-left-right preorder and retains the first
matching shell. It does not skip a malformed numerical payload in favor of a
later witness. It extracts the $\mathsf{LLLR}$ audit, applies the public carrier
grammar, and restores the deleted true bit. At the designated phase the recovered
bitword starts with $10$ and then complete 19-bit one-hot blocks. Restoring the
known selected symbol 18 yields the UT19 event word.

The numerical reader finds the first maximal run of symbol 16. For a halting
BP2 counter value $x$, its length is $4^{x+1}$. The reader tests the power-of-four
form and recovers $x$. The register compiler's output convention is $x=2v+3$,
so the fixed final arithmetic $(x-3)/2$ returns register result $v$. The reader
rejects malformed lengths or an invalid odd-value convention. It never
constructs $4^x$ to search for a match, reruns the source, or searches a trace.
Lemma~\ref{lem:first} makes the first-preorder policy sound even when the first
accepted tree has multiple matching occurrences.

\subsection{Completion of the main proofs}
\begin{proof}[Proof of Theorems~\ref{thm:register} and~\ref{thm:tape}]
The finite-register compiler preserves and reflects halting and gives the
literal BP2 result $2v+3$. The assembled UT19 simulation preserves and reflects
BP2 halting; its nonhalting proof covers every microstep, and its first event
has the stated counter-run grammar. One-hot simulation identifies the first
CTS event and excludes earlier exhaustion. Proposition~\ref{prop:generic}
provides the actual native-S trajectory; Lemmas~\ref{lem:origins} and
\ref{lem:first} identify its earliest accepted sample and every witness payload.
The fixed numerical reader therefore returns $v$ there.

For the converse, suppose no designated CTS event occurs. The compiler's
nonexhaustion certificate then makes every CTS prestate nonempty.
The global stage invariant therefore applies to every finite S prefix and
excludes its event pattern. Contraposition gives a CTS event from an accepted
S sample, hence a source halt by the two source simulations. Halting
absorption and determinism give a unique register result. Equality of first
acceptance indices then yields the value-specific equivalence in
Theorem~\ref{thm:register}. This is the actual composition in
\code{SOnlyUniversality}; source liveness and numerical-success premises of
component lemmas have been discharged by the compiled source.

The concrete tape-to-three-register compiler preserves every finite run
and every represented tape cell. Positive macro durations make its checkpoints cofinal; an internal
register halt is reflected to a source halt by absorption. Composing its
\code{input_halting_iff} with the register result proves
Theorem~\ref{thm:tape}. Its register 0 is exactly the specified finite left-stack
code, giving the stated TM-specific numerical result.
\end{proof}

\section{Effective bounds}\label{sec:resources}

\subsection{Encoding}
For $R$ source registers and $Q$ labels, put $k=R+5Q$, $N_c=4k+1$ and $c=8k+3$.
If $s$ is the sum of the natural inputs and $L$ is the literal BP2 program
length, the actual uniform-helper compiler has the checked bound
\[
 L\leq 2c(2s+9)+N_c(18c+2)+c+4N_c+4.
\]
It chooses memory depth $d=\lfloor\log_2(L+2)\rfloor+2$, so the width satisfies
$W=2^d\leq4L+8$. The compact initializer and all initial counters have size
polynomial in $c,L,W$; large source values are initialized by the compiled
program's execution, rather than by inserting their exponentially expanded
counter representation into the seed. With fixed $P$ and $A$, S generation
has size and construction cost linear in the bitword, up to fixed code
constants.

For an explicit binary description of length $b$, source dimensions are bounded
by the tabulated description size and $s$ is at most singly exponential in $b$.
Consequently the complete syntax encoder has a bound $2^{\operatorname{poly}(b)}$.
The tape compiler has
$\mathrm{states}\cdot\mathrm{symbols}\cdot(5(\mathrm{symbols}+1)+5)+1$ labels.
An $n$-symbol stack code is at most $(\mathrm{symbols}+1)^n-1$, so its binary
length is $O(n\log(\mathrm{symbols}+1))$. These facts preserve a singly
exponential bound for the conventional-tape encoder as well. This is a
source-description bound, independent of whether or when the source halts.

\subsection{Observation and readout}
The fixed bottom-up observer of Lemma~\ref{lem:regular} takes linear time in
occurrence-tree size, with a construction-dependent constant.
For a decoder input of size $m$, witness search visits at most $m$ occurrences
and uses a fixed pattern at each node. The selected audit is a subterm.

The carrier decoder has at most one recursive continuation per layer, always
on a strict descendant, and reaches at most one base cell-spine decoder.
Thus there are at most $m$ charged carrier/cell iterations and at most $m$
decoded bits. A classifier may compare arbitrary continuation subtrees, and
the append-based spine-argument function can take quadratic work on a malformed
long spine. A fixed quadratic allowance per iteration gives a safe
$O((m+1)^3)$ constructor/list bound. The remaining fixed-block parser and
numerical tests have polynomial bit cost; dynamic numerical counts are bounded
by the decoded word length. This establishes a polynomial all-input algorithm,
including malformed terms and rejected payloads.

The proof deliberately separates two levels. \code{SOnlyDecoderBound} checks
strict size decrease, output-length bounds, fuel equivalence, charged recursion
and quadratic spine visits. \code{SOnlyDecoderBoundCurrent.readFuel_eq} binds
the bounded copy to the actual first-preorder decoder on every input. The
remaining local-parser operation estimates and standard bit-cost interpretation
are written arguments. The cubic statement is not a fully instrumented Lean
runtime theorem. Upstream's different row/terminal/bit readers already have a
fully instrumented quartic certificate~\cite{quartic}; comparing their exponents
requires aligning both functions and cost models.

\section{Formal artifact and reproducibility}\label{sec:verification}

The development uses official Lean 4.33.1~\cite{lean} and the generic upstream source at
commit \code{85a867988442fc423279341200f81634a1e65582}. There are 67 project
modules, including eleven additions for the conventional-tape bridge and its
checks, and a 423-module pinned upstream dependency closure. No claim is made
that this dependency replay covers the upstream release's complete
1,239-module headline package.

A clean build of all project modules and an import-only wrapper completed
successfully. Official {\small\texttt{leanchecker \mbox{-}\mbox{-}fresh}} replay then accepted the
aggregate. A separate audit queried 1,359 explicit public declarations and
types; dependencies were restricted to \code{propext}, \code{Quot.sound} and
\code{Classical.choice}. The encoder, finite-label compiler, observer and
current-tree decoder have no choice dependency in their data cones. No new
axiom, unfinished proof or native-decision shortcut is used.

The portable helper independently rebuilt all 490 upstream/project source
modules plus its wrapper in a fresh isolated workspace with no old project or
dependency objects mounted, and then passed fresh kernel replay. It verifies
pinned source/configuration hashes before building, refuses a nonempty output
directory, executes no Lake configuration, and uses only the selected official
library and fresh outputs. Recorded observations are:
\begin{center}
\begin{tabular}{lrr}
\toprule
Check & Wall seconds & Peak sampled RSS\\
\midrule
67 project modules and wrapper & 57.170 & 1.135 GB\\
Aggregate fresh replay & 125.013 & 1.475 GB\\
Full portable source build and replay & 633.802 & 1.516 GB\\
\bottomrule
\end{tabular}
\end{center}
These are measurements of proof checking, not of simulated source computation.
The portable run's replay portion took 136.142 seconds. Independent rehashing
covered 426 pinned source/config files, 424 dependency outputs, all 17,505
installed toolchain files, and 136 final source/output/wrapper files. The
isolated proof runs used read-only proof inputs and denied socket creation.

Positive and negative controls used the same unchanged checker. A valid proof
passed, an ordinary invalid proof of False failed source checking, and an
explicitly injected invalid theorem was rejected by fresh replay for a kernel
type mismatch. Canary objects were excluded from the proof search path.
Fresh replay is an additional pass through the official kernel, not an
independently bootstrapped alternate kernel.

Independent semantic review examined the source compiler, all-microstep halt
reflection, first-event composition and exact output scope. Bounded transcriptions
checked 29,400 assembled transitions and 135,048 tape/stack steps, among other
cases. All 721 current Python test methods passed across a recorded completed
prefix and resumed suffix; the earlier complete 693-method baseline also
passed. These tests detect transcription defects and support reproduction;
they do not replace the quantified formal statements.

The source packet includes the pinned upstream sources and MIT notice, exact
module identities, theorem types, declaration dependencies, command records,
all relevant logs, and the portable build helper. The principal result manifest
is \code{results/formal_turing_universality.json}; the register result remains
\code{results/formal_universality.json}. The main source identity is recorded
in Appendix~\ref{app:reproduce}.

\section{Scope and further questions}

The result concerns one explicitly prescribed finite-controller selected path.
It establishes neither strategy-independent simulation nor universal computation
through S normal-form termination. The finite control, root reset, local
observations, read-only navigation, native-only mutations and invocation bound
expose exactly what the reduction strategy contributes. A bounded next-step
routine alone would not settle every possible interpretation of auxiliary
computational power; the stated controller model is part of the theorem.

The compact phase count provides a concrete alternative instantiation, not a
minimality theorem or a measured global speedup. Sharper controller-state and
native-contraction comparisons, a fully instrumented cost semantics for the
current numerical reader, and a canonical tape-output normalization theorem
are separate potential extensions. No unproved conjecture is used in the
construction.

The S-combinator challenge announcement expressly discusses selected paths and
observation during nonterminating evolution~\cite{challenge}. Whether a
particular construction satisfies the committee's intended admissibility and
originality criteria is an adjudication question, not a theorem of the formal
model. The contribution and attribution above are mathematical claims; this
paper asserts no prize award or priority over the earlier public universality
construction.

\clearpage
\appendix
\section{Theorem entry points and reproduction}\label{app:reproduce}

The shortest semantic entry points are:
\begin{itemize}
\item \code{SOnlyTuringUniversality.halting_iff_event}: arbitrary finite-alphabet
conventional tape source and finite input, with the fixed S event.
\item \code{SOnlyTuringUniversality.selector_executes} and
\code{every_step_native}: the actual chosen sequence and native steps.
\item \code{SOnlyUniversality.result_iff_first_event}: exact arbitrary-register
numerical output at the first accepted tree.
\item \code{SOnlyTuringUniversality.left_stack_result_iff_first_event}: the
explicit finite-left-stack convention.
\item \code{SOnlyGlobalFirstEvent.firstSourceEvent_current_read}: sharp event
creation and common first-event audit.
\item \code{SOnlyTapeCompiler.input_halting_iff}: conventional tape source to
the concrete three-register machine.
\end{itemize}

With Python 3.10 or newer, the official Lean 4.33.1 distribution and the included
pinned dependency source directory, the reproducible command is:
\begin{verbatim}
python tools/replay_formal.py --target turing \
  --lean-bin /absolute/path/to/lean-4.33.1-linux/bin \
  --upstream-source /absolute/path/to/upstream \
  --build-dir /absolute/path/to/new-empty-replay
\end{verbatim}
The optional \code{--check-only} verifies identities and dependency order but
performs no proof check. The helper is a direct offline build procedure, not
itself an OS sandbox. Its selected source manifests are pinned by hash. It
does not automatically repeat the separately recorded declaration audits or
negative controls.

The main theorem source \code{formal/SOnlyTuringUniversality.lean} has SHA-256
\begin{quote}\small\nolinkurl{9a8b1e653f5e5f0c206a4be305deb24be51e7e1f7ae93e9e9a200862d02056a3}.\end{quote}
The 67-module result manifest has SHA-256
\begin{quote}\small\nolinkurl{588778633d2b5a26288a4b77ed870dd2730a282e0ff67fadfd64bcb4a913916a}.\end{quote}
The recorded official Linux x86-64 Lean distribution archive has SHA-256
\begin{quote}\small\nolinkurl{0376ac87487246b40dd077268c097e701f552e94c6d020d2373b50c7444fa22f}.\end{quote}

\begin{thebibliography}{99}
\raggedright
\bibitem{strawberry}
Cinematic Strawberry. \emph{Pure S Is Computationally Universal Under a Fixed
Root-Restarted Finite Controller}. Manuscript dated 14 September 2026;
Alexander Filin, project coordinator. Pinned repository revision
\code{85a867988442fc423279341200f81634a1e65582}.
\href{https://github.com/cstrawberry/predictive-universe/blob/85a867988442fc423279341200f81634a1e65582/docs/paper/related/pure_s_universality/paper.md}{Manuscript and supporting MIT-licensed formalization}.

\bibitem{headline}
Cinematic Strawberry. \code{PureSFormal.RootResetHeadline} and
\code{PureSFormal.RootResetChallenge}. In the formalization of~\cite{strawberry}.
\href{https://github.com/cstrawberry/predictive-universe/blob/85a867988442fc423279341200f81634a1e65582/docs/paper/related/pure_s_universality/formalization/PureSFormal/RootResetHeadline.lean}{Pinned headline source}.

\bibitem{infinite}
Cinematic Strawberry. \code{DeterministicTapeInfiniteTape}.
Conventional integer-indexed Boolean tape bridge in~\cite{strawberry}.
\href{https://github.com/cstrawberry/predictive-universe/blob/85a867988442fc423279341200f81634a1e65582/docs/paper/related/pure_s_universality/formalization/PureSFormal/Computation/DeterministicTapeInfiniteTape.lean}{Pinned source}.

\bibitem{prioroutput}
Cinematic Strawberry. \code{DeterministicTapeRootResetOutput} and
\code{RootResetRegularMarkerLanguage}. Exact current-tree observations and
regular halting in~\cite{strawberry}.
\href{https://github.com/cstrawberry/predictive-universe/blob/85a867988442fc423279341200f81634a1e65582/docs/paper/related/pure_s_universality/formalization/PureSFormal/Computation/DeterministicTapeRootResetOutput.lean}{Output source};
\href{https://github.com/cstrawberry/predictive-universe/blob/85a867988442fc423279341200f81634a1e65582/docs/paper/related/pure_s_universality/formalization/PureSFormal/Research/RootResetRegularMarkerLanguage.lean}{regular-language source}.

\bibitem{generic}
Cinematic Strawberry. \code{RootResetFiniteAllInputsTraceAgreement}.
Generic finite-CTS realization in~\cite{strawberry}.
\href{https://github.com/cstrawberry/predictive-universe/blob/85a867988442fc423279341200f81634a1e65582/docs/paper/related/pure_s_universality/formalization/PureSFormal/Research/RootResetFiniteAllInputsTraceAgreement.lean}{Pinned source}.

\bibitem{quartic}
Cinematic Strawberry. \code{DeterministicTapeOutputQuartic}.
Instrumented all-input output resource certificate in~\cite{strawberry}.
\href{https://github.com/cstrawberry/predictive-universe/blob/85a867988442fc423279341200f81634a1e65582/docs/paper/related/pure_s_universality/formalization/PureSFormal/Computation/DeterministicTapeOutputQuartic.lean}{Pinned source}.

\bibitem{ut19}
ais523. \emph{UT19}. Esolang Wiki, revision 154545. CC0 production table
and source construction.
\url{https://esolangs.org/w/index.php?title=UT19&oldid=154545}.

\bibitem{bp2}
ais523. \emph{Brainpocalypse II}. Esolang Wiki, revision 172507.
\url{https://esolangs.org/w/index.php?title=Brainpocalypse_II&oldid=172507}.

\bibitem{waterfall}
ais523. \emph{The Waterfall Model}. Esolang Wiki, revision 193038;
the account credits zzo38's independent earlier model.
\url{https://esolangs.org/w/index.php?title=The_Waterfall_Model&oldid=193038}.

\bibitem{amnesiac}
ais523. \emph{The Amnesiac From Minsk}. Esolang Wiki, revision 173298.
\url{https://esolangs.org/w/index.php?title=The_Amnesiac_From_Minsk&oldid=173298}.

\bibitem{lean}
Lean FRO. \emph{Lean 4.33.1}. Official toolchain release.
\url{https://github.com/leanprover/lean4/releases/tag/v4.33.1}.

\bibitem{challenge}
Stephen Wolfram. \emph{1920--2020 and a \$20,000 Prize: Announcing the
S Combinator Challenge}. 8 June 2021.
\href{https://writings.stephenwolfram.com/2021/06/1920-2020-and-a-20000-prize-announcing-the-s-combinator-challenge/}{Official announcement}.
\end{thebibliography}

\end{document}
